% Part: normal-modal-logic
% Chapter: syntax-and-semantics
% Section: normal-modal-logics

\documentclass[../../../include/open-logic-section]{subfiles}

\begin{document}

\olfileid{nml}{syn}{nor}

\olsection{सामान्य मोडल तर्कशास्त्रे}

मोडल तर्कशास्त्र म्हणजे योग्य अटी पूर्ण करणारा
!!{formula}s चा संच असे मानता येते. काही मोडल
तर्कशास्त्रे विशेष महत्त्वाची आहेत: त्यांची अर्थलक्षणे
आणि उपयोग रंजक आहेत, आणि विशेषतः त्यांचा
सैद्धान्तिक अभ्यास चांगला झालेला आहे. यांनाच
सामान्य मोडल तर्कशास्त्रे म्हणतात.

\begin{defn}\ollabel{defn:modal-logic}
\emph{मोडल तर्कशास्त्र}~$\Log L$ म्हणजे मूलभूत
मोडल भाषेतील !!{formula}s चा असा संच की
\begin{enumerate}
\item त्यात सर्व विधानात्मक सर्वत:सत्य सूत्रे असतात;
\item तो एकरूप आदेशनाखाली बंद असतो; आणि
\item तो मोडस पोनेन्सखाली बंद असतो; म्हणजे $!A$ आणि
  $(!A \lif !B) \in \Log L$ असतील, तर
  $!B \in \Log L$ सुद्धा असते.
\end{enumerate}
\end{defn}

\begin{ex}
फारसे रंजक नसले तरी ही व्याख्या पूर्ण करणारे,
म्हणून मोडल तर्कशास्त्रे असणारे, !!{formula}s
चे काही संच असे आहेत:
\begin{enumerate}
\item सर्व सर्वत:सत्य सूत्रांचा संच;
\item सर्व !!{formula}s चा संच (विसंगत मोडल तर्कशास्त्र).
\end{enumerate}
\end{ex}

\begin{defn}\label{defn:normal-modal-logic}
मोडल तर्कशास्त्र~$\Log L$ हे तेव्हा आणि केवळ
तेव्हाच \emph{सामान्य} असते, जेव्हा
\begin{enumerate}
\item त्यात पुढील दोन्ही असतात:
\begin{align}
\tag{K} \Box(p \lif q) & \lif (\Box p \lif \Box q) \\
\tag{Dual} \Diamond p & \liff \lnot \Box \lnot p
\end{align}
\item आणि ते अनिवार्यीकरणाखाली बंद असते; म्हणजे
  $!A \in \Log L$ असेल, तर
  $\Box !A \in \Log L$ सुद्धा असते.
\end{enumerate}
\end{defn}

मोडल तर्कशास्त्रे परिभाषित करण्याचे दोन प्रमुख
मार्ग आपण पाहू. पहिला सिद्धता-उपपत्तीय: विशिष्ट
!!{derivation} पद्धतींमध्ये निष्पन्न करता येणाऱ्या
!!{formula}s चा संच. दुसरा प्रतिमान-उपपत्तीय:
प्रतिमानांच्या विशिष्ट वर्गांत वैध असणाऱ्या
!!{formula}s चे संच. साध्या विधानात्मक तर्कशास्त्रात
(किंवा प्रथम-क्रम तर्कशास्त्रात) हाच प्रकार दिसतो.
तेथेही !!{formula}s चा संच दोन प्रकारे ठरवता येतो:
सर्व सत्यतामूल्य-नियुक्त्यांत सत्य असलेल्या
सर्वत:सत्य सूत्रांचा संच म्हणून, किंवा एखाद्या
!!{derivation} पद्धतीतील सर्व !!{derivable}
!!{formula}s चा संच म्हणून. सामान्य मोडल
तर्कशास्त्रांतही संबंधीय प्रतिमाने आणि स्वयंसिद्धकी
(हिल्बर्ट-प्रकारच्या) सिद्धता पद्धती वापरून हे
दुहेरी वर्णन शक्य होते.
\end{document}
